\addtocontents{toc}{\protect\Needspace{3\baselineskip}}
\originalchapter{8.821 课程详讲：尺度、视界与温度}
\label{ch:holo-geometry-full}
本章对应第 2--3 讲: Planck 尺度、经典黑洞的有关性质、Rindler 近视界、共形因果图以及欧氏正则性求温度. Einstein 方程和 Schwarzschild 解本身的求解采用已有 GR 卷作为前置, 本章围绕后续量子态与热力学所需的坐标、能量和因果联系展开, 不重写 GR 全卷. 默认 $c=1$, 涉及物理温度时保留 $\hbar,\kb$.

\originalsection{量子长度与引力长度为何交会}
以 $G_N,\hbar,c$ 构造出的三个单位是
\[
m_P=\sqrt{\frac{\hbar c}{G_N}},\qquad
\ell_P=\sqrt{\frac{\hbar G_N}{c^3}},\qquad
t_P=\frac{\ell_P}c.
\]
量纲可逐项核对: $[G_N]=L^3M^{-1}T^{-2}$、$[\hbar]=ML^2T^{-1}$、$[c]=LT^{-1}$. 这些是此处使用的\textbf{未约化} Planck 单位; 部分引力文献把质量定义为 $\sqrt{\hbar c/(8\pi G_N)}$, 二者不可不说明就互换. 数量级为 $m_Pc^2\simeq1.22\times10^{19}\,\mathrm{GeV}$、$\ell_P\simeq1.62\times10^{-35}\,\mathrm m$、$t_P\simeq5.39\times10^{-44}\,\mathrm s$. 原讲义的质量指数或文本提取易出错, 本卷以定义式和 SI 常数核对.

对质量 $m$, 约化 Compton 长度 $r_C=\hbar/(mc)$ 描述相对论量子定位的尺度, 引力长度则由 Schwarzschild 半径 $r_s=2G_Nm/c^2$ 给出. 比值是
\[
\frac{r_s}{r_C}=2\left(\frac m{m_P}\right)^2.
\]
当 $m\ll m_P$, 试图把粒子定位至 $r_s$ 以前量子效应已经重要, 不能用一个小经典 Schwarzschild 黑洞替代通常粒子. 当 $m\gg m_P$, $r_s\gg r_C$, 宏观黑洞可在小曲率条件下作半经典计算. 当 $m$ 与 $m_P$ 同量级, 两者交会, 一套同时含量子与引力效应的描述不可避免.

在高能散射尺度 $E$, 无量纲引力强度约为 $\alpha_G\sim G_NE^2/(\hbar c^5)$. 对静质量可写 $\alpha_G=G_Nm^2/(\hbar c)=(m/m_P)^2$. 电磁相互作用强度 $\alpha\simeq1/137$ 在这种比较中不随粒子静质量平方变化. 电子的引力强度因此极小, 但这不意味着宏观黑洞没有可观测量子效应: 它的温度很小而熵很大, 两者由视界组织.

\begin{example}
把局域探测的长度估计为 $\delta x\sim\hbar/\delta p$, 探针能量取超相对论 $E\sim c\delta p$, 其引力半径则是 $r_s\sim2G_N\delta p/c^3$. 要求探针不在所测尺度内形成强引力区, 粗略要 $\delta x\gtrsim r_s$, 两式给 $\delta x^2\gtrsim2\ell_P^2$. 这是最小可探测长度的量级\textbf{启发}, 不是对所有测量装置、量子态与量子引力理论成立的严格最短长度定理; 系数也不由这种估计可靠确定.
\end{example}
经典引力对应忽略量子效应; 固定曲背景 QFT 保留 $\hbar$ 而忽略场对几何的回馈; 半经典引力则先量子化物质, 并在适当范围用其能动张量期望值修正几何. 三种近似对同一问题的适用条件不同, 后面的温度推导主要处于后两者.

\originalsection{Schwarzschild 视界的红移与有限曲率}
已有 GR 解为
\[
\dd s^2=-f(r)\dd t^2+\frac{\dd r^2}{f(r)}+r^2\dd\Omega_2^2,
\qquad f(r)=1-\frac{r_s}r,\quad r_s=2G_NM.
\]
此处 $t$ 在无穷远归一化为静止观察者固有时. 视界 $r=r_s$ 处 $g_{tt}=0$、$g_{rr}$ 发散, 但曲率不变量
\[
R_{\mu\nu\rho\sigma}R^{\mu\nu\rho\sigma}
=\frac{12r_s^2}{r^6}
\]
在视界有限, 在 $r=0$ 发散. 视界的坐标问题与中心的真实曲率奇点由此区分. 完整永恒解对时间反演对称, 真实由坍缩形成的黑洞则有不同的过去区域; 使用两侧永恒黑洞解释热纠缠时要保留这一区别.

在 $r>r_s$ 静止的观察者有 $u^\mu=f^{-1/2}(\partial_t)^\mu$, 固有时满足 $\dd\tau=\sqrt f\dd t$. 粒子的 Killing 能量 $E_\infty=-p_t$ 守恒, 局部能量则为
\[
E_{\rm local}=-p_\mu u^\mu=\frac{E_\infty}{\sqrt f},
\qquad E_\infty=\sqrt f\,E_{\rm local}.
\]
固定局部能量的过程移近视界, 无穷远观察到的能量趋零. 这是无限红移, 不是说自由落体的局部时钟本身失效. 严格停在视界需要无限加速度, 不存在该处正常的静止时向观察者.

这个加速度可直接求出. $\Gamma^r{}_{tt}=ff'/2$, 因而 $a^r=\Gamma^r{}_{tt}(u^t)^2=f'/2$, 模长为
\[
a(r)=\sqrt{g_{rr}(a^r)^2}=\frac{|f'|}{2\sqrt f}.
\]
乘红移因子后取视界极限, 定义表面引力
\[
\kappa=\lim_{r\to r_s^+}\sqrt f\,a(r)
=\frac{f'(r_s)}2=\frac1{2r_s}=\frac1{4G_NM}.
\]
它不是视界处某个静止观察者可直接读出的有限固有加速度; 其有限性依靠无穷远归一化的红移. 若改变 $t$ 的常数归一化, $\kappa$ 也随 Killing 生成元归一化改变.

径向自由落体守恒量 $e=f\dot t$ 和归一化给 $\dot r^2=e^2-f$. 视界附近 $\dot r\to-e$ 有限, 而 $\dd t/\dd r\sim-1/f$ 对数发散. 因此跨视界固有时有限、Schwarzschild 坐标时无限, 两者是同一轨迹的不同参数描述. PS1.3 给出余下寿命和光信号的完整考核.

\originalsection{Rindler 坐标从近视界展开得来}
定义沿定时间径向片从视界量起的固有距离 $\rho$, 近非极端视界有
\[
\dd\rho=\frac{\dd r}{\sqrt f},\qquad
\rho\simeq2\sqrt{\frac{r-r_s}{f'(r_s)}},\qquad
f\simeq\kappa^2\rho^2.
\]
于是径向时间部分是 $-\kappa^2\rho^2\dd t^2+\dd\rho^2$. 令无量纲 boost 时间 $\eta=\kappa t$, 则
\[
\dd s^2\simeq-\rho^2\dd\eta^2+\dd\rho^2+r_s^2\dd\Omega_2^2.
\]
变换 $T=\rho\sinh\eta$、$X=\rho\cosh\eta$ 后, 前两项成为 $-\dd T^2+\dd X^2$. 正 $\rho$ 只覆盖 $X>|T|$ 的右楔, 不是全部 Minkowski 平面. $\rho\to0$、有限 $\eta$ 只到原点; 要到 $X=T>0$ 的未来视界, 还需同时 $\eta\to\infty$、$\rho\ee^\eta$ 保持有限. 这解释了为什么一个“坐标边界”包含整条零曲线.

\begin{center}
\begin{tikzpicture}[scale=0.92,>=Latex]
\fill[structurecolor!5] (0,0)--(2.2,2.2)--(2.2,-2.2)--cycle;
\draw[->] (-2.5,0)--(2.5,0) node[right] {$X$};
\draw[->] (0,-2.4)--(0,2.4) node[above] {$T$};
\draw[dashed] (-2.2,-2.2)--(2.2,2.2);
\draw[dashed] (-2.2,2.2)--(2.2,-2.2);
\draw[structurecolor,thick,domain=-1.12:1.12,samples=70,variable=\s]
plot ({cosh(\s)},{sinh(\s)});
\node[font=\small] at (1.65,0.35) {$\rho=1$};
\node[font=\small] at (1.5,-1.9) {右 Rindler 楔};
\node[font=\small] at (-1.7,0.5) {左楔};
\end{tikzpicture}
\end{center}
图中的双曲线由确定公式绘制, 虚线是零方向. 对固定 $\rho$ 轨迹, $\tau=\rho\eta$, 四速度为 $(\cosh\eta,\sinh\eta)$, 四加速度模为 $a=1/\rho$. 因此静止黑洞观察者近视界对应匀加速观察者; 自由落体近视界则近似惯性轨迹. 温度的观察者依赖将在下一章由量子态解释.

\originalsection{共形因果图保留哪些信息}
二维 Minkowski 取 $U=T-X,V=T+X$, 度规为 $-\dd U\dd V$. 令 $u=\arctan U,v=\arctan V$, 两者范围都为 $(-\pi/2,\pi/2)$, 有
\[
\dd s^2=-\frac{\dd u\dd v}{\cos^2u\cos^2v}.
\]
乘正共形因子 $\Omega^2=\cos^2u\cos^2v$, 得 $-\dd u\dd v$. 再用纵轴 $u+v$、横轴 $v-u$, 无穷大时空被压进有限菱形, 零曲线仍倾斜 $45^\circ$. 因为正共形缩放保持光锥, 图能判断谁可向谁发信号; 它并不保存固有距离或时间数值.

在该图, $i^\pm$ 表示定空间位置趋向未来/过去无限固有时间的端点, $i^0$ 是空间无穷远, $\mathcal I^\pm$ 是出入射零射线的无穷远边界. 对四维球对称问题画二维径向图时, 每个内点实际代表一个二球, 不能把图上的面积当成四维空间体积.

为看 Schwarzschild 延拓, 用 $\dd r_*/\dd r=1/f$, 即
\[
r_*=r+r_s\log\left|\frac r{r_s}-1\right|,
\quad U_K=-\ee^{-(t-r_*)/(2r_s)},\quad
V_K=\ee^{(t+r_*)/(2r_s)}.
\]
外部满足 $-U_KV_K=(r/r_s-1)\ee^{r/r_s}$, 度规径向部分为
\[
\dd s^2_{(2)}=-\frac{4r_s^3}{r}\ee^{-r/r_s}\dd U_K\dd V_K.
\]
其系数在 $r=r_s$ 有限, 延拓穿过 $U_K=0$ 或 $V_K=0$ 后得到两个外部、未来黑洞区和过去白洞区. 再用有限范围坐标与共形因子压缩, 得下面的因果示意图.

\begin{center}
\begin{tikzpicture}[scale=0.83,>=Latex]
\fill[structurecolor!6] (0,0)--(-1.7,1.7)--(1.7,1.7)--cycle;
\draw[thick] (-3.4,0)--(-1.7,1.7)--(1.7,1.7)--(3.4,0)--(1.7,-1.7)--(-1.7,-1.7)--cycle;
\draw[dashed] (-1.7,-1.7)--(1.7,1.7);
\draw[dashed] (-1.7,1.7)--(1.7,-1.7);
\draw[ultra thick] (-1.7,1.7)--(1.7,1.7);
\draw[ultra thick] (-1.7,-1.7)--(1.7,-1.7);
\node at (1.8,0) {右外部};\node at (-1.8,0) {左外部};
\node at (0,1) {未来黑洞};\node at (0,-1) {过去白洞};
\node[above,font=\small] at (0,1.7) {$r=0$ 奇点};
\node[below,font=\small] at (0,-1.7) {$r=0$ 奇点};
\node[right] at (3.4,0) {$i^0$};\node[left] at (-3.4,0) {$i^0$};
\node[above right,font=\small] at (2.6,0.8) {$\mathcal I^+$};
\node[below right,font=\small] at (2.6,-0.8) {$\mathcal I^-$};
\end{tikzpicture}
\end{center}
虚线为视界, 加粗水平边界为时空的空间型奇点. 未来指向曲线从右外部跨未来视界后, 只能到上方奇点, 无法再到右侧 $\mathcal I^+$. 过去白洞可以发出信号但不是一个坍缩黑洞必须具备的真实区域. 图是因果结构示意, 不是宣称把 Schwarzschild 全度规等距画进欧氏平面.

\originalsection{无锥尖的欧氏时间周期}
有限温度量子力学的配分函数是 $Z(\beta)=\operatorname{Tr}\ee^{-\beta H}$. 路径积分将它写为欧氏时间首尾粘合的传播, 周期是 $\hbar\beta$; 玻色场周期, 费米场反周期. 所以欧氏时间\textbf{周期结构}可揭示平衡温度, 但还需选定对应的量子态.

将近视界 $t=-\ii t_E$, 径向时间度规变成
\[
\dd s_E^2=\dd\rho^2+\kappa^2\rho^2\dd t_E^2.
\]
这是极坐标平面, 角度为 $\theta=\kappa t_E$. 若角度周期不是 $2\pi$, 原点会有锥形缺口或重叠. 正则性要求 $t_E$ 周期 $2\pi/\kappa$, 与 $\hbar\beta$ 比较, 得
\[
T_\infty=\frac{\hbar\kappa}{2\pi\kb}
=\frac{\hbar}{8\pi\kb G_NM}.
\]
恢复 $c$ 后为 $T_\infty=\hbar c^3/(8\pi\kb G_NM)$. 温度的参考时钟必须与前面 $t$、$\kappa$ 的归一化一致.

局部固有欧氏时间周期缩短为 $\sqrt f\hbar\beta$, 故 Tolman 关系是
\[
T_{\rm local}(r)=\frac{T_\infty}{\sqrt{f(r)}},\qquad
T_{\rm local}\sqrt{-g_{tt}}=\text{常数}.
\]
\textbf{局部温度和局部能量都除以红移因子}; 不能一边说趋近视界发散, 一边误写成乘 $\sqrt f$. 这是核查原公式和文本提取的简单一致性检查. 近视界 $\sqrt f\simeq\kappa\rho$, 因而 $T_{\rm local}\simeq\hbar/(2\pi\kb\rho)=\hbar a/(2\pi\kb)$, 即 Unruh 温度, 此处 $c=1$.

\begin{example}
静止探测器从 $f=1/4$ 移到 $f=1/100$, 局部平衡温度分别为 $2T_\infty$ 和 $10T_\infty$. 它读到的差别不是无穷远热浴变了温度, 而是采用了不同红移的固有时钟. 自由落体探测器并不沿这些静止轨迹, 因此不能从静止探测器温度发散就推断正则视界上所有观察者都遇到无限热墙.
\end{example}
PS2.1 将同一正则性步骤推广至不同 $f,h$ 的非极端视界. 极端视界的二阶零点改变了局部几何, 其极限和量子态需另外处理; PS2.2 给出对应检验. 下一章进一步说明欧氏正则性选出的热态从何而来.
